\documentclass{article}
\usepackage[utf8]{inputenc}
\usepackage{graphicx}
\usepackage[headheight=2pt,headsep=0pt]{geometry}

\title{Contrat de location du local}
\author{Cercle Informatique de l'Université Libre de Bruxelles}
\date{Cerkinfo A.S.B.L}

\begin{document}

\begin{figure}
	\includegraphics[scale=0.03]{images/ci.png}
	\hfill
	\includegraphics[scale=0.4]{images/ulb.png}
  \label{fig:ci}
\end{figure}

\maketitle

\section*{Bailleur}
Désigné le \underline{BAILLEUR} de ce contrat :

\begin{center}
    \textbf{Cerkinfo A.S.B.L}\\
    \textbf{Boulevard du Triomphe - CP 206 - 1050 Ixelles}
\end{center}
Représenté par :

    \textbf{\textit{Nom :}}\\
    
    \textbf{\textit{Prénom :}}\\
    
    \textbf{\textit{Poste dans le comité Cerkinfo :}}

\section*{Locataire}
Désigné le \underline{LOCATAIRE} de ce contrat : \\

    \textbf{\textit{Cercle ACE :}}\\

Représenté par :

    \textbf{\textit{Nom :}}\\
    
    \textbf{\textit{Prénom :}}\\
    
    \textbf{\textit{Poste dans le comité de cercle :}}
    
\newpage

\section*{ARTICLE 1 - Durée}
Le \underline{LOCATAIRE} loue la salle à partir du ...../...../..... à .....H..... jusqu’au ...../...../..... à .....H..... 
Afin que l’état des lieux d’entrée puisse être dressé, il s’engage à se présenter le jour du début de la location. À la fin de location, le \underline{LOCATAIRE} devra restituer la salle à l’heure prévue. Il devra rester le temps nécessaire pour permettre l’établissement de l’état des lieux de sortie. Il s’engage à rendre les lieux dans un état de propreté identique à celui qu’il aura constaté lors du début de la location.

\section*{ARTICLE 2 - Local}
Le présent contrat concerne la salle dénommée \textbf{Cerkinfo} et située : 
\begin{center}
    \textbf{Campus La Plaine ULB - Boulevard du Triomphe - 1050 Ixelles}\\
    \textbf{Local \underline{P.OF.082}}
\end{center}

\section*{ARTICLE 3 - Charte Horaire}
Le \textbf{Cerkinfo} est soumis à la Charte Horaire des cercles de l'Association des Cercles Etudiants. En conséquence, le \underline{LOCATAIRE} s'engage à respecter les horaires de fermeture du local, à savoir 0h30 fin de la soirée (coupure de la musique et plus de service de boisson), et 1h fin du nettoyage et fermeture définitive. En cas de non respect de cette charte, le \underline{LOCATAIRE} assumera toute responsabilité vis-à-vis des autorités de l'ULB.

\section*{ARTICLE 4 - Équipement}
Le matériel mis à disposition doit être rendu propre et en bon état de fonctionnement. Il fera l’objet d’un inventaire lors des états des lieux d’entrée et de sortie de la salle.

\section*{ARTICLE 5 - Obligations}
Le \underline{BAILLEUR} est tenu de mettre le local à la disposition du \underline{LOCATAIRE} à la date et à l’heure convenue pour le début de la location.


\section*{ARTICLE 6 - Versement}
Le \underline{LOCATAIRE} s’engage à :

\begin{minipage}{0.6\textwidth}
\begin{itemize}
    \item Payer le dépôt de garantie, en cash lors de la signature du présent contrat ou effectuer un versement sur le compte \textbf{BE22 3632 1093 6947} avant la signature du présent contrat (cfr. ARTICLE 7)
    \item Payer le nombre de fûts commandés, en cash lors de la signature du présent contrat ou effectuer un versement sur le compte \textbf{BE22 3632 1093 6947} avant la signature du présent contrat (cfr. ARTICLE 8)
\end{itemize}
\end{minipage}
\begin{minipage}{0.4\textwidth}\raggedleft
\includegraphics[scale=0.04]{images/QRCodeBancontact.png}
\end{minipage}\\\\
Il est précisé que le présent contrat sera \textbf{\underline{résilié}} de plein droit si une somme quelconque n’est pas payée à son échéance.

\newpage

\section*{ARTICLE 7 - Caution}
Le \underline{LOCATAIRE} s’engage à verser une caution de ............. €, payable en cash à la signature du contrat ou sur le compte bancaire \textbf{\underline{au préalable}}. Cette somme sera restituée après l’état des lieux de sortie. En cas de dommages observés, le montant des \underline{réparations}, du \underline{nettoyage} ou du \underline{remplacement du matériel} sera déduit du dépôt de garantie. Si le montant du préjudice est supérieur au montant de la caution, le locataire s’engage à rembourser les frais supplémentaires sous \textbf{6 jours ouvrables} après constatation des dégâts.

\section*{ARTICLE 8 - Fûts}
Le \underline{BAILLEUR} met à disposition du \underline{LOCATAIRE} ......... fût(s) de bière Jupiler au prix de ........ €/fût. Le \underline{BAILLEUR} s'engage a rembourser tout fût non consommé si celui-ci n'as pas été décacheté et entamé.\\\\
Une majoration de \textbf{15€/fut} sera ajouté à tout fut entamé au delà de ceux mit à disposition.

\section*{ARTICLE 9 - Contentieux}
Tout litige relatif à l’exécution du présent contrat sera soumis à défaut d’accord amiable aux tribunaux dont dépend le siège social du \underline{BAILLEUR}.

\section*{ARTICLE 10 - Sous-location}
Il est interdit au \underline{LOCATAIRE} de consentir une quelconque location de la salle.

\section*{ARTICLE 11 - Accès}
Le Président de Cerkinfo A.S.B.L ou son représentant ou le délégué ci-dessus indiqué, aura \underline{libre accès} à la salle durant toute la durée de la location.

\section*{ARTICLE 12 - Nettoyage}
Le \underline{LOCATAIRE} s’engage à rendre le matériel et les locaux aussi propres que dans l’état initial et à vider les poubelles. En cas de manquement à cette clause, le \underline{LOCATAIRE} s’engage à verser un supplément de \textbf{50€}.

\section*{ARTICLE 13 - Loi}
Le \underline{LOCATAIRE} s’engage à ne commettre ou à ne laisser dans la salle, aucun acte contraire à la loi.

\section*{ARTICLE 14 - Fumer}
Le Cercle est un espace non fumeur.  En conséquence, il est strictement prohibé de consommer des cigarettes ou des joints au sein de cet espace. Le \underline{LOCATAIRE} accepte de prendre l'entière responsabilité des conséquences résultant de toute activation de l'alarme incendie dans le cadre de cette interdiction.

\section{ARTICLE 15 - Nuisances sonores}
Le \underline{LOCATAIRE} s'engage à respecter les riverains en limitants les nuisances sonores à l'extérieur du local. Cela comprends la limitation du volume sonore des enceintes et la limitation des cris et des bruits de la part des participants à l'extérieur. Dans la mesure du possible, le \underline{LOCATAIRE} s'engage à fermer les portes pour limiter les nuisances sonores. En cas d'intervention de la police suite à une plainte des riverains, ou en cas d'intervention d'un·e agent de sécurité du campus, le \underline{LOCATAIRE} s'engage à assumer la pleine responsabilité auprès des autorités de l'ULB, en se reconnaissant comme unique responsable de l'événement.

\section*{État des lieux}
\begin{center}
\end{center}
\begin{center}
\begin{tabular}{ |c|c|c| } 
 \hline
 \textbf{Local} & \textbf{Status Entrée} & \textbf{Status Sortie} \\
 \hline
 Murs &  &  \\
 \hline
 Canapés &  &  \\
 \hline
 Armoires &  &  \\
 \hline
 Sound System &  &  \\
 \hline
 Armoire sono & & \\
 \hline
 Bar &  &  \\
 \hline
 Pompe &  &  \\
 \hline
 Frigo &  &  \\
 \hline
 Extincteur et trousse de secours &  &  \\
 \hline
  &  &  \\ 
  \hline
  &  &  \\ 
  \hline
  &  &  \\ 
 \hline
 \textbf{Décompte} & \textbf{Nbre Entrée} & \textbf{Nbre Sortie} \\
 \hline
 Bombonne Gaz Cacheté &  &  \\
 \hline
 Bombonne Gaz Vide/Entamé &  &  \\
 \hline
 Fut Jupiler Cacheté &  &  \\
 \hline
 Fut Jupiler Entamé &  &  \\
 \hline
 Fut Jupiler Vide &  &  \\
 \hline
  &  &  \\
 \hline
  &  &  \\
  \hline
  &  &  \\ 
 \hline
\end{tabular}
\end{center}

\section*{Signature}
Fait à .......................... \\\\
Le ......./......../........ \\\\
Faire précéder la signature de la mention manuscrite \textbf{"Lu et approuvé"}\\\\

\begin{minipage}{0.5\textwidth}
\underline{Le LOCATAIRE}
\end{minipage}
\begin{minipage}{0.5\textwidth}\raggedleft
\underline{Représentant Cerkinfo A.S.B.L}
\end{minipage}

\end{document}
